\chapter{Analisi Difensiva}
\label{ch:difensiva}

Il Capitolo~\ref{ch:scenari} ha documentato tre scenari di call stack spoofing
producendo, per ciascuno, le evidenze osservate da tre strumenti distinti.
Questo capitolo cambia prospettiva: per ogni scenario identifica la firma rilevabile,
il meccanismo di difesa appropriato e il livello di privilegio richiesto per applicarlo.
La sezione~\ref{sec:mapping} illustra firma e meccanismo di difesa per ciascuna tecnica;
la sezione~\ref{sec:edr-kernel} descrive la logica kernel-mode con cui un EDR
copre l'intero spazio delle tecniche analizzate.

% ── 4.1 Mapping Difensivo ─────────────────────────────────────────────────────

\section{Meccanismi di Difesa per Scenario}
\label{sec:mapping}

Per ciascuna tecnica di spoofing analizzata esiste una firma osservabile e un
meccanismo di difesa corrispondente. La distinzione fondamentale è tra
\emph{rilevazione} e \emph{blocco}: un componente user-mode può osservare
anomalie in modo asincrono, ma al momento del rilevamento l'operazione
sospetta potrebbe essere già completata. Il blocco affidabile richiede che il
controllo avvenga in modo sincrono nel contesto del thread, in kernel-mode.

\subsection{Single-Frame Spoofing}

Single-Frame Spoofing modifica unicamente il LR salvato (\texttt{[x29+8]})
nel frame di \texttt{InjectExplorer}, sostituendo l'indirizzo di ritorno reale
con quello di un gadget \texttt{ntdll}. Il puntatore di frame \texttt{x29}
non viene modificato: continua a puntare al frame reale del chiamante in
\texttt{stack\_spoof.exe}.

La firma è nella catena \texttt{x29} stessa. Un walker che percorra la catena
\texttt{x29}/\texttt{x30} risale da \texttt{InjectExplorer} al frame
precedente tramite \texttt{[x29]} e trova \texttt{stack\_spoof.exe} come
chiamante reale: l'attribuzione al codice offensivo non è nascosta dalla
tecnica. Qualsiasi walker che segua \texttt{x29} anziché fidarsi unicamente
del LR rilevato è sufficiente per identificare l'origine reale. Il blocco
sincrono richiede tuttavia un livello di privilegio superiore.

\subsection{Stack Pivot}

La firma prodotta dallo scenario Stack Pivot è deterministica: al momento della
syscall, SP contiene l'indirizzo del fake stack allocato con
\texttt{VirtualAlloc}, esterno all'intervallo
\texttt{[TEB.StackLimit, TEB.StackBase]}. Nessun programma legittimo opera con
SP fuori dai bounds TEB durante una syscall.

Il meccanismo di difesa è la lettura di \texttt{TrapFrame.Sp} in un callback
kernel. Al momento della transizione user-mode/kernel, il processore ARM64
salva il valore esatto di SP nella struttura \texttt{\_KTRAP\_FRAME}, allocata
nello stack kernel del thread e non modificabile da user-mode. Un callback
registrato tramite \texttt{PsSetCreateThreadNotifyRoutine} viene eseguito nel
contesto del thread chiamante, prima che il kernel completi l'operazione: in
quel momento è possibile leggere \texttt{TrapFrame.Sp}, confrontarlo con i
limiti TEB, e restituire un codice di errore che annulla l'intera operazione.
Il controllo è deterministico e privo di falsi negativi.

% ── 4.2 EDR Kernel-Mode ──────────────────────────────────────────────────────

\section{EDR in kernel-mode: rilevamento e blocco}
\label{sec:edr-kernel}

Gli strumenti analizzati nel Capitolo~\ref{ch:scenari} condividono un limite
strutturale: operano nello stesso spazio di indirizzi del processo che
ispezionano. Le strutture che leggono (la catena \texttt{x29}/\texttt{x30},
i record \texttt{.pdata}, i limiti TEB) sono le stesse strutture che il
codice offensivo ha manipolato per costruire lo stack falsificato. Un walker
user-mode non può distinguere tra struttura autentica e struttura artefatta,
perché entrambe sono dati dello stesso spazio di memoria.

A questo si aggiunge un secondo limite, indipendente dal primo: l'ispezione
avviene in modo asincrono rispetto all'esecuzione del thread bersaglio. Quando
il rilevamento si concretizza, l'operazione sospetta è già completata. Si può
osservare, non prevenire.

Un EDR di classe enterprise affronta entrambi i problemi con un approccio
qualitativamente diverso: opera in kernel-mode e si aggancia al flusso di
esecuzione in modo sincrono. Questi due aspetti non sono equivalenti e vanno
tenuti separati.

\subsection{Spazio di esecuzione, fonti di osservazione e sincronia}

Il driver EDR opera in kernel space e legge strutture che nessun codice
user-mode può raggiungere o modificare: lo stato interno del thread, le
tabelle di unwinding, il \texttt{\_KTRAP\_FRAME} allocato sullo stack kernel
alla transizione syscall (Sezione~\ref{sec:kernel-stack}). I walker user-mode
leggono invece strutture nello stesso spazio di indirizzi del processo
offensivo: la catena \texttt{x29}/\texttt{x30}, i record \texttt{.pdata}, i
limiti TEB sono tutti dati che l'attaccante ha già manipolato prima che il
walker li osservi.

Il campo \texttt{TrapFrame.Sp} fornisce il valore autoritativo di SP al
momento della syscall. Il walker kernel-mode percorre la catena
\texttt{x29}/\texttt{x30} con lo stesso meccanismo di
\texttt{CaptureStackBackTrace}: legge alternativamente il frame pointer
salvato e il return address a ogni livello. La differenza è nel punto di
partenza: il primo frame pointer non viene letto dall'x29 live del thread,
ma da \texttt{KTRAP\_FRAME.X29}, catturato dall'hardware alla transizione
syscall e non modificabile da user-mode.

Il driver si registra come \emph{callback di notifica}: il kernel lo invoca
in modo sincrono nel contesto del thread offensivo, prima che l'operazione
abbia effetto. Nei tre scenari analizzati il punto di intercettazione
critico è \texttt{NtCreateThreadEx}, la syscall sottostante a
\texttt{CreateRemoteThread}: è in questo momento che lo stack falsificato è
attivo e il driver effettua i controlli per individuare se l'operazione è
benigna o malevola.

Se il controllo rileva un'anomalia, può restituire un codice di errore che il
kernel propaga al chiamante: l'operazione viene annullata, non semplicemente
registrata. Il thread offensivo non può distinguere questo scenario da un
fallimento ordinario della syscall. Non riceve notifica dell'intervento; vede
solo che l'operazione non è andata a buon fine.

\subsection{Applicazione agli scenari analizzati}

\paragraph{Stack Pivot.}
La firma è deterministica e rilevabile tramite il trap frame. Al momento della
syscall che crea il thread remoto, SP contiene l'indirizzo del fake stack
allocato con \texttt{VirtualAlloc}, esterno all'intervallo
\texttt{[TEB.StackLimit, TEB.StackBase]}. Il controllo su SP è privo di falsi
negativi, poiché nessun programma legittimo invoca una syscall con SP fuori
dai bounds TEB, e si traduce in un blocco sincrono prima della creazione del
thread.

\paragraph{Single-Frame Spoofing.}
La firma è rilevabile tramite la catena \texttt{x29}/\texttt{x30}. Single-Frame
Spoofing sostituisce il LR salvato di \texttt{InjectExplorer} con un gadget
\texttt{ntdll}, ma non modifica il puntatore di frame \texttt{x29}: seguendo
\texttt{[X29]} a partire dal frame di \texttt{InjectExplorer}, il walker
raggiunge il frame reale del chiamante in \texttt{stack\_spoof.exe}. Il driver
legge \texttt{KTRAP\_FRAME.X29} come punto di partenza autoritativo, non
modificabile da user-mode, e percorre la catena fino a rivelare l'attribuzione
reale. Il controllo è deterministico: \texttt{stack\_spoof.exe} è sempre
presente nella catena \texttt{x29}, indipendentemente dalla manipolazione del LR.
